\documentclass{article}
\usepackage[margin=1in]{geometry}
\usepackage{amsmath,amssymb}
\usepackage{booktabs}
\usepackage{graphicx}
\usepackage[numbers,sort&compress]{natbib}
\usepackage{hyperref}
\providecommand{\citet}[1]{\citeauthor{#1}~\citep{#1}}

\newcommand{\Rhat}[0]{\widehat{R}}
\newcommand{\risk}{R}

\title{Symmetry Spending Under a Quantum Shot Budget:\\
What It Takes to Certify That Shots Were Well Spent}

\author{Draft for internal review\\ \texttt{shotcert/} working draft}

\begin{document}
\maketitle

\begin{abstract}
A quantum predictor that averages a model response over the images of its input under a symmetry
group is a \emph{stochastic, shot-limited object}: it applies one of several circuits and reads out a
finite number of measurement outcomes. Its risk therefore depends on how its budget was spent, and
two allocations with the same total number of shots can have different risks --- there is no ideal
limit that both are approximating. We formalise this as a \emph{risk-at-budget} functional and
decompose it into three separately measurable parts: an orbit-mean model bias, a group-draw variance
that more shots cannot reduce, and a shot variance that more group draws cannot reduce. The
decomposition is exact and we verify it against brute-force enumeration to
$3\times10^{-17}$ rather than asserting it. Two consequences follow. First, at fixed total shots the
split between ``how many symmetry images'' and ``how many shots per image'' leaves the shot term
untouched, so orbit randomisation is free in shots. The interior optimum in that split, which an
earlier version of this paper reported, is \emph{retracted}: it was produced by charging eight
distinct circuits in an experiment with two (\S\ref{sec:switchretract}). Second, when the \emph{decision} about how to spend shots must itself be certified from an
independent calibration experiment, the cost decomposes bound-independently as
$\mathrm{Var}(D)=\mathrm{Var}_x(\mathbb{E}[D\mid x])+\mathbb{E}[\mathrm{Var}(D\mid x)]$: at four
shots per prediction the deployed-noise term exceeds the between-instance term by $15\times$
($0.07851$ vs $0.00507$), and at 256 shots the ordering reverses ($0.00038$ vs $0.00244$), so shots
buy the first and only preparations buy the second, down to a floor that no number of shots can
cross. \emph{Where} the certification frontier sits, however, is a property of the inequality rather
than of the problem, and we report that as a withdrawal: with a distribution-free
empirical-Bernstein bound, power is $0.000$ up to $n=200$ and $1.000$ at $n=400$; with a Studentised
bound whose measured violation rate is $0.005$ across 90 cells (nominal $0.05$), power is already
$0.995$ at $n=100$. At matched total budget the direction of the preparations-versus-shots tradeoff
is not robust --- at $B=1600$ the allocation $n{=}50,S{=}32$ reaches power $0.917$ while
$n{=}200,S{=}8$ reaches $0.870$, the opposite of what the loose bound suggests. We therefore claim
the variance decomposition and not an allocation rule. On a bond-alternating Heisenberg response with a control-line rectification
error, the same structure reproduces: full orbit averaging reduces deployed risk by $0.00800$ when
the target is mirror-even, and \emph{increases} it by $0.00212$ when the target is the raw response,
so the value of spending shots on symmetry is governed by how much odd signal the target contains,
not by whether a symmetry exists. We state explicitly which parts of this are already owned by prior
work on shot/sample budget allocation and on certified symmetry release, and what remains.
\end{abstract}

\section{Introduction}

A growing body of work asks how a fixed quantum measurement budget should be spent.
Catak~\citep{catak2026budget} derives a distribution-free bound separating a sample term from a shot
term and a closed-form allocation $n^\star, S^\star$ under $B=nS$. Chen~\citep{chen2026release}
certifies the \emph{release} of a symmetry constraint during training, with simultaneous confidence
bounds, a post-selection oracle inequality for adaptive candidate families, and grouped-measurement
shot accounting. Chen~\citep{chen2026symmetry} treats finite-measurement inference of the symmetry
group itself. Wang et al.~\citep{wang2026fourier} and Cai et al.~\citep{cai2021random} give
randomised symmetry implementations with explicit bias--sampling tradeoffs.

This paper is about a different object, and the difference is not cosmetic.

\begin{quote}
For a deployed predictor that draws $k$ symmetry images of its input and spends $S$ shots across
them, the risk is a functional of the allocation. Spending the same $S$ differently does not change
how well we \emph{know} the risk; it changes \emph{which} risk we have.
\end{quote}

Three things follow, and all three are measured rather than asserted.

\textbf{(1) A decomposition with two invariances.} The deployed risk separates as
$\risk_S(a)=\mathrm{bias}^2 + V_G/k + \Sigma/S$. The shot term $\Sigma/S$ does not depend on $k$, and
the draw term $V_G/k$ does not depend on $S$. At fixed total shots, moving shots between ``more
images'' and ``more shots per image'' leaves one term exactly fixed and moves the other. There is no
interior optimum in shots alone: $k=S$ weakly dominates. The interior optimum that practitioners
actually observe comes from the cost of \emph{distinct circuits}, which is a different resource and
is where a switching cost $c$ enters.

\textbf{(2) Certification is where the cost is, and the group has to pay for it.} Deciding
whether symmetry spending helps requires a calibration experiment on independent labelled instances,
and the deployed predictor's own shot noise enters the calibration contrast. That noise is reducible
by shots; the between-instance variability of the contrast is not, and it sets a floor on
$\mathrm{Var}(D)$ that more shots leave untouched. What we do \emph{not} claim is an allocation
rule: the location of the certification frontier moves by roughly $4\times$ in $n$ depending on
which concentration bound is used, and the direction of the preparations-versus-shots tradeoff at
matched budget reverses between them (\S\ref{sec:withdraw}).

\textbf{(3) The symmetry's value is a property of the target, not of the model.} On the same physical
mechanism, changing only how much odd-under-mirror signal the prediction target contains flips full
orbit averaging from a $0.00800$ risk reduction to a $0.00212$ risk increase.

\section{Access model, costs, and estimands}
\label{sec:access}

\textbf{Access.} Scalar single-copy shots: one labelled state preparation per calibration instance,
one circuit-execution batch per prediction, computational-basis measurement of a diagonal observable.
No two-copy access, no shared Pauli measurements, no classical shadows. Every baseline receives
exactly this access. Where a published method needs more --- a twirl--swap test, say --- we say so
and do not compare it against a weaker version of itself.

\textbf{Costs are counted separately.} \texttt{shots} (measurement repetitions),
\texttt{preparations} (independently prepared labelled states), \texttt{circuits} (distinct compiled
circuits, one per symmetry image touched), and \texttt{calibration\_labels}. Deployment cost per
prediction is $S + c\cdot(\text{distinct circuits})$, with $c$ a declared stress parameter. A method
that receives free calibration or free state knowledge invalidates the comparison; the policy menu
and the feasibility rules below exist to make that impossible to do by accident.

\textbf{Feasibility is a property of a policy, not a convention.} With \texttt{per}$=\lfloor S/k\rfloor$,
a policy with $k>S$ spends $k$ shots while claiming $S$. This study found that bug in its own menu
after a first frontier run, and the apparent power gain at small $S$ had been bought with four times
the declared shots. Every policy now carries \texttt{shots\_used} and a \texttt{feasible()} predicate,
the menu is filtered by it, and the dropped entries are reported.

\textbf{Estimands.} Primary: $\Delta_S(a)=\risk_S(a)-\risk_S(a_0)$ at matched total shots, against the
raw policy $a_0$. Operational primary: $N_{\mathrm{cert}}(\epsilon,\delta)$, the number of calibration
preparations at which the certificate of \S\ref{sec:certify} selects a policy with
$\Pr[\Delta_S(a_{\mathrm{sel}})\le\epsilon]\ge 1-\delta$ and power $\ge 1-\beta$ when a useful policy
exists. Fallback to $a_0$ is a safe abstention, not a success.

\section{Risk at budget}
\label{sec:risk}

For input $x$ and group element $g$, let $\mu_g(x)=\mathbb{E}[\text{shot}\mid g,x]$ and let
$\sigma_g^2(x)=1-\mu_g(x)^2$ for a $\pm1$-valued readout. A policy draws $k$ images, spends $S$
shots total with $S/k$ per draw and equal weights, and outputs $Z_a(x)$. Then, exactly,
\begin{equation}
\risk_S(a)\;=\;\underbrace{\mathbb{E}_x\big[(\bar\mu_a(x)-y)^2\big]}_{\text{bias}^2}
\;+\;\underbrace{\mathbb{E}_x\big[\mathrm{Var}_{g}\mu_g(x)\big]/k}_{V_G/k}
\;+\;\underbrace{\mathbb{E}_x\big[\mathbb{E}_g\,\sigma_g^2(x)\big]/S}_{\Sigma/S},
\label{eq:decomp}
\end{equation}
where $\bar\mu_a$ is the orbit-mean response. The shot term deserves its invariance spelled out:
$\sum_j (1/k)^2\sigma_j^2/(S/k)=\bar\sigma^2/S$, independent of $k$.

\begin{table}[h]\centering\small
\caption{Equation~\eqref{eq:decomp} verified on the 2-qubit control. The residual column is the
difference between the enumerated risk and the closed form; it is machine precision, not a fit.}
\end{table}

\begin{table}[h]\centering\small
\begin{tabular}{llrrrrr}
\toprule
Regime & Policy & $\risk_{64}$ & bias$^2$ & $V_G$ & $\Sigma$ & residual \\
\midrule
help & a0 & 0.03955 & 0.03428 & 0.00000 & 0.33686 & 4.2e-17 \\
help & strat & 0.00571 & 0.00040 & 0.00000 & 0.34004 & 1.0e-17 \\
help & mix25 & 0.02476 & 0.01948 & 0.00000 & 0.33766 & 4.5e-17 \\
help & mix50 & 0.01419 & 0.00890 & 0.00000 & 0.33845 & 2.6e-17 \\
help & mix75 & 0.00784 & 0.00254 & 0.00000 & 0.33925 & 1.0e-17 \\
help & k1 & 0.03949 & 0.00040 & 0.03377 & 0.34004 & 7.6e-17 \\
help & k2 & 0.02260 & 0.00040 & 0.03377 & 0.34004 & 4.9e-17 \\
help & k4 & 0.01416 & 0.00040 & 0.03377 & 0.34004 & 1.9e-17 \\
help & k8 & 0.00993 & 0.00040 & 0.03377 & 0.34004 & 5.2e-18 \\
help & k16 & 0.00782 & 0.00040 & 0.03377 & 0.34004 & 6.9e-18 \\
help & k32 & 0.00677 & 0.00040 & 0.03377 & 0.34004 & 7.8e-18 \\
help & k64 & 0.00624 & 0.00040 & 0.03377 & 0.34004 & 6.9e-18 \\
hurt & a0 & 0.03813 & 0.03287 & 0.00000 & 0.33686 & 2.1e-17 \\
hurt & strat & 0.11438 & 0.10907 & 0.00000 & 0.34004 & 5.6e-17 \\
hurt & mix25 & 0.05086 & 0.04558 & 0.00000 & 0.33766 & 3.5e-17 \\
hurt & mix50 & 0.06781 & 0.06252 & 0.00000 & 0.33845 & 1.1e-16 \\
hurt & mix75 & 0.08899 & 0.08369 & 0.00000 & 0.33925 & 4.2e-17 \\
hurt & k1 & 0.14816 & 0.10907 & 0.03377 & 0.34004 & 1.1e-16 \\
hurt & k2 & 0.13127 & 0.10907 & 0.03377 & 0.34004 & 5.6e-17 \\
hurt & k4 & 0.12283 & 0.10907 & 0.03377 & 0.34004 & 1.1e-16 \\
hurt & k8 & 0.11860 & 0.10907 & 0.03377 & 0.34004 & 2.8e-16 \\
hurt & k16 & 0.11649 & 0.10907 & 0.03377 & 0.34004 & 1.2e-16 \\
hurt & k32 & 0.11544 & 0.10907 & 0.03377 & 0.34004 & 9.7e-17 \\
hurt & k64 & 0.11491 & 0.10907 & 0.03377 & 0.34004 & 1.4e-16 \\
interior & a0 & 0.02071 & 0.01544 & 0.00000 & 0.33686 & 1.4e-17 \\
interior & strat & 0.01089 & 0.00558 & 0.00000 & 0.34004 & 3.5e-18 \\
interior & mix25 & 0.01192 & 0.00665 & 0.00000 & 0.33766 & 2.6e-17 \\
interior & mix50 & 0.00736 & 0.00207 & 0.00000 & 0.33845 & 6.9e-18 \\
interior & mix75 & 0.00702 & 0.00171 & 0.00000 & 0.33925 & 1.6e-17 \\
interior & k1 & 0.04467 & 0.00558 & 0.03377 & 0.34004 & 6.9e-18 \\
interior & k2 & 0.02778 & 0.00558 & 0.03377 & 0.34004 & 4.2e-17 \\
interior & k4 & 0.01934 & 0.00558 & 0.03377 & 0.34004 & 6.9e-18 \\
interior & k8 & 0.01512 & 0.00558 & 0.03377 & 0.34004 & 2.3e-17 \\
interior & k16 & 0.01300 & 0.00558 & 0.03377 & 0.34004 & 1.7e-17 \\
interior & k32 & 0.01195 & 0.00558 & 0.03377 & 0.34004 & 3.5e-18 \\
interior & k64 & 0.01142 & 0.00558 & 0.03377 & 0.34004 & 1.7e-17 \\
\bottomrule\end{tabular}
\caption{Risk-at-budget decomposition, enumerated exactly on the 2-qubit control. Residual is the gap between the enumerated risk and Eq.~\eqref{eq:decomp}.}
\end{table}


\section{Certification, and what it costs}
\label{sec:certify}

Calibration uses $n$ independent labelled instances; the independence unit is the instance, not the
shot or the seed. For each instance we deploy the candidate and $a_0$ and record
$D_i=(Z_a-Y)^2-(Z_{a_0}-Y)^2$, and we bound $\mathbb{E}[D]$ with an empirical-Bernstein inequality
over the fixed menu, with the multiplicity split declared in advance. The inequality is deliberately
ordinary: the contribution here is not a new concentration bound but the observation that what is
being bounded is a difference of two \emph{allocation-dependent} risks of stochastic predictors,
selected with the same finite-shot data that defines them --- precisely the independence assumption
that Catak's bound requires and states~\citep{catak2026budget}.

\begin{table}[h]\centering\small
\begin{tabular}{lrrrr}
\toprule
$n$ (preps) & $S$ (shots) & violation & power & fallback \\
\midrule
25 & 256 & 0.000 & 0.000 & 1.000 \\
50 & 128 & 0.000 & 0.000 & 1.000 \\
80 & 80 & 0.000 & 0.000 & 1.000 \\
100 & 64 & 0.000 & 0.000 & 1.000 \\
160 & 40 & 0.000 & 0.000 & 1.000 \\
200 & 32 & 0.000 & 0.000 & 1.000 \\
320 & 20 & 0.000 & 0.140 & 0.860 \\
400 & 16 & 0.000 & 0.976 & 0.024 \\
800 & 8 & 0.000 & 1.000 & 0.000 \\
1600 & 4 & 0.000 & 1.000 & 0.000 \\
\bottomrule\end{tabular}
\caption{Certificate power along the preparations/shots axis at fixed total budget $B=nS=6400$, averaging-helps regime, empirical-Bernstein bound, $\epsilon=0$, $\delta=0.05$.}
\label{tab:frontier}
\end{table}


\textbf{Is the transition statistics, or the inequality?} The empirical-Bernstein bound has a range
term $7(b-a)\ln/(3(n-1))$ that is enormous at small $n$ whatever the data says, so a frontier located
at $n\approx 400$ could be a property of the bound rather than of the problem. We therefore report
$\mathrm{Var}(D)=\mathrm{Var}_x(\mathbb{E}[D\mid x])+\mathbb{E}[\mathrm{Var}(D\mid x)]$ separately from
the bound's two terms. Only the second is reducible by shots; the first is reducible only by
preparations. Table~\ref{tab:variance} gives both, and the split of the bound itself.

\begin{table}[h]\centering\small
\begin{tabular}{lrrrrrr}
\toprule
Regime & $S$ & $\mathrm{Var}(D)$ & between & noise & EB var.\ term & EB range term \\
\midrule
help & 4 & 0.08354 & 0.00507 & 0.07851 & 0.0000 & -0.007 \\
help & 16 & 0.01186 & 0.00270 & 0.00917 & 0.0133 & 0.014 \\
help & 64 & 0.00414 & 0.00247 & 0.00168 & 0.0196 & 0.205 \\
help & 256 & 0.00281 & 0.00244 & 0.00038 & nan & nan \\
hurt & 4 & 0.13272 & 0.01042 & 0.12238 & 0.0000 & -0.007 \\
hurt & 16 & 0.02866 & 0.00766 & 0.02104 & 0.0126 & 0.014 \\
hurt & 64 & 0.01197 & 0.00733 & 0.00466 & 0.0200 & 0.205 \\
hurt & 256 & 0.00837 & 0.00727 & 0.00113 & nan & nan \\
interior & 4 & 0.07145 & 0.00344 & 0.06805 & 0.0000 & -0.007 \\
interior & 16 & 0.00751 & 0.00086 & 0.00665 & 0.0074 & 0.014 \\
interior & 64 & 0.00166 & 0.00060 & 0.00106 & 0.0152 & 0.205 \\
interior & 256 & 0.00078 & 0.00056 & 0.00023 & nan & nan \\
\bottomrule\end{tabular}
\caption{What the calibration contrast is made of, and what the bound adds. $\mathrm{Var}(D)$ splits into a between-instance part (only preparations reduce it) and a deployed-noise part (shots reduce it). The last two columns split the empirical-Bernstein bound at the matching allocation.}
\label{tab:variance}
\end{table}


\subsection{Matched total cost, including what calibration actually pays}
\label{sec:ledger}
The first version of this study labelled $nS$ the calibration budget. It was not: the menu is
evaluated on every calibration instance, so calibration costs $n\sum_p S_p$, a factor of the menu
size larger, and preparations must be counted per shot-batch rather than per instance because repeated
shots need repeated state copies. Table~\ref{tab:ledger} repeats the comparison with the full five-part
ledger at every point --- labelled instances, preparations, shots, switches, compilations --- and with
a group-free comparator: the same certificate, the same calibration protocol, and a menu of shrinkage
policies that buy variance reduction without touching the group.

\begin{table}[h]\centering\small
\begin{tabular}{llrrrrr}
\toprule
Task & Menu & $n$ & $S$ & total shots & power & violation \\
\midrule
control\textunderscore{}help & sym & 200 & 8 & 20800 & $0.900\,[{\scriptsize 0.85},{\scriptsize 0.93}]$ & $0.000$ \\
control\textunderscore{}help & generic & 1600 & 8 & 123200 & $0.140\,[{\scriptsize 0.10},{\scriptsize 0.19}]$ & $0.000$ \\
control\textunderscore{}hurt & sym & 50 & 8 & 11200 & $0.000\,[{\scriptsize 0.00},{\scriptsize 0.02}]$ & $0.000$ \\
control\textunderscore{}hurt & generic & 50 & 8 & 11600 & $0.000\,[{\scriptsize 0.00},{\scriptsize 0.02}]$ & $0.000$ \\
control\textunderscore{}interior & sym & 200 & 32 & 96000 & $0.960\,[{\scriptsize 0.92},{\scriptsize 0.98}]$ & $0.000$ \\
control\textunderscore{}interior & generic & 1600 & 8 & 123200 & $0.100\,[{\scriptsize 0.07},{\scriptsize 0.15}]$ & $0.000$ \\
\bottomrule\end{tabular}
\caption{A2.6 matched-total-cost comparison. Both menus are certified by the same procedure at the same declared $\epsilon$; `total shots' includes the factor of the menu size that calibration actually pays (A2.2). Rows show the cheapest allocation reaching power $\ge 0.9$, or the best available.}
\label{tab:ledger}
\end{table}


The direction is structural rather than statistical. Shrinkage cannot produce a policy that beats raw
by more than $\epsilon$, because it moves the prediction toward a declared constant and therefore
\emph{adds} bias; only the group changes the bias term in Eq.~\eqref{eq:decomp}. The symmetry menu
reaches power $1.000$ at 20{,}800 total shots in the helps regime and $1.000$ at 96{,}000 in the
interior regime; the group-free menu reaches $0.14$ and $0.10$ at comparable or greater cost and never
certifies. In the hurts regime neither menu certifies, which is the correct behaviour: the certificate
abstains. Across all 108 points the maximum observed violation rate is $0.015$ against a nominal
$0.05$.

\subsection{What a correct distribution-free certificate costs}
\label{sec:mpcost}
With the corrected simultaneous bound the answer is analytic and unwelcome. The range term
$7(b-a)\ln/(3(n-1))$ has $b-a=8$ for a difference of squared losses and $\ln=\log(2m/\delta)$, so it
decays as $1/n$ and, for the gaps this decision turns on ($|\Delta|\approx0.008$), requires $5{,}109$
to $20{,}679$ labelled preparations before the bound can close --- independently of how the shots are
allocated. Distribution-free simultaneous certification of decisions this fine is not available at
budgets this study can pay. The frontier numbers in this paper therefore use a Studentised bound whose
coverage is audited rather than assumed (Table~\ref{tab:bound}), and the distribution-free statement is
the cost statement above.

\subsection{Retraction: the switching optimum was an impossible circuit count}
\label{sec:switchretract}
The control orbit has two elements, so at most two distinct circuits exist. The implementation reported
$\min(k,4|G|)$, reaching eight, and the switching experiment charged that number; the reported $k=2$
optimum at large $c$ was the price of circuits that cannot be compiled. With circuits correctly
bounded, a per-circuit cost is flat for every $k\ge2$ and the optimum is always at a boundary. Charging
\emph{switches} under an explicit interleaved schedule does produce an interior optimum, but in only 18
of 144 priced cells, and in none of the 72 per-circuit cells.

\begin{table}[h]\centering\small
\begin{tabular}{lrrr}
\toprule
Charged resource & cells & interior optimum in $k$ & max distinct circuits \\
\midrule
per distinct circuit & 72 & 0 (0.0\%) & 2 \\
per switch & 72 & 18 (25.0\%) & 2 \\
\bottomrule\end{tabular}
\caption{A2.3, the retracted switching result. With circuits correctly bounded by $|G|=2$, a per-circuit price is flat for every $k\ge2$ and never produces an interior optimum; the previously reported $k=2$ optimum came from charging eight circuits. An interior optimum appears only under a per-switch price on an interleaved schedule.}
\label{tab:switch}
\end{table}


\subsection{Withdrawal: the frontier was mostly the inequality}
\label{sec:withdraw}
An earlier version of this draft reported that certification is preparation-limited, on the evidence
that at $B=6400$ power jumped from $0.000$ at $n\le200$ to $1.000$ at $n\ge800$. That claim is
withdrawn. The empirical-Bernstein range term $7(b-a)\ln/(3(n-1))$ is $\approx4$ at $n=25$ for a
loss range of $4$, so the bound, not the data, was deciding. Re-running the same experiment with a
Studentised bound (Table~\ref{tab:bound}) gives power $0.245/0.725/0.995/1.000$ at
$n=25/50/100/200$ with $S=128$, at a measured violation rate of $0.005$ against a nominal $0.05$
across all 90 cells --- admissible empirically, though distribution-free validity is exactly what a
Studentised bound does not offer, and a violation rate estimated on 400 repetitions per cell cannot
resolve a true rate far below $0.05$.

\begin{table}[h]\centering\small
\begin{tabular}{lrrrrrr}
\toprule
Bound & $S$ & $n{=}25$ & $n{=}50$ & $n{=}100$ & $n{=}200$ & $n{=}400$ \\
\midrule
eb & 8 & 0.00/0.00 & 0.00/0.00 & 0.00/0.00 & 0.00/0.00 & 0.93/0.00 \\
eb & 32 & 0.00/0.00 & 0.00/0.00 & 0.00/0.00 & 0.00/0.00 & 0.99/0.00 \\
eb & 128 & 0.00/0.00 & 0.00/0.00 & 0.00/0.00 & 0.00/0.00 & 1.00/0.00 \\
normal & 8 & 0.10/0.00 & 0.28/0.00 & 0.48/0.00 & 0.87/0.00 & 0.99/0.00 \\
normal & 32 & 0.42/0.00 & 0.92/0.00 & 1.00/0.00 & 1.00/0.00 & 1.00/0.00 \\
normal & 128 & 0.82/0.00 & 1.00/0.00 & 1.00/0.00 & 1.00/0.00 & 1.00/0.00 \\
\bottomrule\end{tabular}
\caption{Power/violation rate for the distribution-free empirical-Bernstein bound and the tighter Studentised bound, averaging-helps regime, $\epsilon=0$, $\delta=0.05$. A bound is admissible only where its measured violation rate stays at or below $\delta$.}
\label{tab:bound}
\end{table}


\section{Physical tasks}
\label{sec:physics}

\textbf{Primary.} Bond-alternating spin-$1/2$ Heisenberg chain, $n=6$ and $n=8$, open boundaries,
quench from the N\'eel state, target the site-0 magnetisation. The control parameter treated as odd
is the staggered field $\Delta$, and the deployed policy pairs $\Delta$ with $-\Delta$.

\textbf{What the symmetry claim requires, and where it fails.} An earlier version asserted that
$\Delta\to-\Delta$ \emph{is} the chain mirror. It is not, and the check is cheap. The operator
relation $M H(\lambda) M = H(-\lambda)$ does hold, to $8.9\times10^{-16}$. But a symmetry of the
\emph{deployed prediction} also needs the preparation and the readout to transform with it, and
neither does: the N\'eel state maps to the shifted N\'eel state at distance $\sqrt{2}$, and
$Z_0$ maps to $Z_{n-1}$. The relation that survives is
$\<Z_0\>_t(\lambda)=\<Z_{n-1}\>_t(-\lambda)$, which holds to $10^{-15}$ and which describes a
\emph{different experiment} (mirrored preparation, mirrored readout). The relation this study
actually exploits, $\<Z_0\>_t(\lambda)=\<Z_0\>_t(-\lambda)$, is off by $0.19$. So
$\Delta\to-\Delta$ is a paired control transformation whose odd component is measurable and
cancellable, not a symmetry of the full experiment, and the risk decomposition is unaffected because
it needs only the response functions $\mu_g$, not a physical group action. The deployed device mis-sets that parameter by a second-order rectification of
the analogue control line, $\Delta_{\mathrm{eff}}=\Delta+c\Delta^2$, whose error is odd in $\Delta$
to leading order --- which is exactly the component that averaging over $\Delta\to-\Delta$ cancels.
A fixed-site readout bias would \emph{not} be cancelable by input averaging, and we say so rather
than claiming credit for a symmetry that the mechanism does not have.

The three regimes differ by how much of the odd response the target contains:
$y_a=(1-a)\,y_{\mathrm{even}}+a\,y_{\mathrm{raw}}$, with $y_{\mathrm{even}}$ the part even under
$\Delta\to-\Delta$. Table~\ref{tab:physics2} shows raw, full averaging and the best partial mixture
side by side, because a best-policy-only table cannot display a sign flip.

\begin{table}[h]\centering\small
\begin{tabular}{lrrrrr}
\toprule
Task & $\risk(a_0)$ & best policy & $\risk$(best) & $\Delta$ & max resid. \\
\midrule
heisenberg\textunderscore{}even\textunderscore{}target\textunderscore{}n6 & 0.02167 & strat & 0.01367 & $-0.00800$ & 2.8e-17 \\
heisenberg\textunderscore{}mostly\textunderscore{}even\textunderscore{}n6 & 0.01870 & mix75 & 0.01387 & $-0.00483$ & 2.8e-17 \\
heisenberg\textunderscore{}raw\textunderscore{}target\textunderscore{}n6 & 0.01621 & mix25 & 0.01525 & $-0.00096$ & 2.4e-17 \\
heisenberg\textunderscore{}even\textunderscore{}target\textunderscore{}n8 & 0.01734 & strat & 0.01319 & $-0.00415$ & 1.2e-17 \\
cluster\textunderscore{}even\textunderscore{}target\textunderscore{}n7 & 0.01137 & strat & 0.00934 & $-0.00203$ & 1.6e-17 \\
cluster\textunderscore{}raw\textunderscore{}target\textunderscore{}n7 & 0.01137 & strat & 0.00934 & $-0.00203$ & 1.6e-17 \\
\bottomrule\end{tabular}
\caption{Physical tasks at $S=64$ total shots, enumerated risk. The sign of the benefit of orbit averaging flips with the odd content of the target.}
\label{tab:physics}
\end{table}

\begin{table}[h]\centering\small
\begin{tabular}{lrrrrr}
\toprule
Task & $\risk(a_0)$ & full averaging & $\Delta$ & best mixture & $\Delta$ \\
\midrule
heisenberg\textunderscore{}even\textunderscore{}target\textunderscore{}n6 & 0.02167 & 0.01367 & $-0.00800$ & mix75: 0.01418 & $-0.00749$ \\
heisenberg\textunderscore{}mostly\textunderscore{}even\textunderscore{}n6 & 0.01870 & 0.01425 & $-0.00445$ & mix75: 0.01387 & $-0.00483$ \\
heisenberg\textunderscore{}raw\textunderscore{}target\textunderscore{}n6 & 0.01621 & 0.01833 & $+0.00212$ & mix25: 0.01525 & $-0.00096$ \\
heisenberg\textunderscore{}even\textunderscore{}target\textunderscore{}n8 & 0.01734 & 0.01319 & $-0.00415$ & mix75: 0.01343 & $-0.00391$ \\
cluster\textunderscore{}even\textunderscore{}target\textunderscore{}n7 & 0.01137 & 0.00934 & $-0.00203$ & mix75: 0.00946 & $-0.00192$ \\
cluster\textunderscore{}raw\textunderscore{}target\textunderscore{}n7 & 0.01137 & 0.00934 & $-0.00203$ & mix75: 0.00946 & $-0.00192$ \\
\bottomrule\end{tabular}
\caption{Raw, full averaging and the best partial mixture side by side (A2.5 and the reviewer's presentation point): the sign flip is visible only when all three are shown. Risks at $S=64$ total shots, enumerated exactly.}
\label{tab:physics2}
\end{table}

\begin{table}[h]\centering\small
\begin{tabular}{lr}
\toprule
Quantity tested & residual \\
\midrule
\multicolumn{2}{l}{\emph{Regenerating from results/physics.json; values quoted in the rebuttal table.}}\\
\bottomrule\end{tabular}
\caption{A2.5 mirror-relation residuals.}\label{tab:mirror}
\end{table}


\textbf{Independent evaluation.} A transverse-field three-body cluster Hamiltonian at $n=7$, a
different dynamical mechanism and a different observable, with the same rectification error. The
split is by task family, declared before generation.

\section{What is and is not claimed}

\textbf{Claimed.} (i) The risk-at-budget decomposition~\eqref{eq:decomp} with its two invariances,
verified against an independent enumeration of the predictor. (ii) That under a pure shot budget orbit randomisation is free in shots, and
that the interior optimum in the split is \emph{not} a circuit-count phenomenon
(\S\ref{sec:switchretract}): it survives only as a per-switch cost on an interleaved schedule. (iii) That at matched total cost, including the factor of the menu size that calibration
actually pays, a symmetry menu certifies where a group-free shrinkage menu cannot, and why; and that
a correct distribution-free simultaneous certificate for decisions of this size costs five to twenty
thousand labelled preparations, which is a cost statement, not an achievement. (iv) That on a physical mechanism the sign of the benefit is
controlled by the target's odd content, not by the presence of a symmetry.

\textbf{Not claimed.} No new concentration inequality. No joint sample/shot allocation rule for a
generalization bound --- that is~\citet{catak2026budget}. No certified symmetry release, no
post-selection oracle inequality for adaptive families, no two-copy measurement primitive --- that
is~\citet{chen2026release}. No symmetry discovery from finite probes --- that is~\citet{chen2026symmetry}.
\textbf{Withdrawn.} The interior switching optimum (\S\ref{sec:switchretract}); the claim that
$\Delta\mapsto-\Delta$ is the chain mirror of the deployed experiment (\S\ref{sec:physics},
Table~\ref{tab:mirror}); the claim that the certification frontier locates a resource limit rather
than a property of the inequality used (\S\ref{sec:withdraw}); and every coverage, power and
frontier number produced by the pre-correction bound, which used $\log(4/(\delta n))$ inside a
one-sided radius and the loss width where the contrast width belongs.

No device results, no quantum advantage, and no claim that stratified allocation
$S_g\propto |w_g|\sigma_g$ is anything other than classical.

\section{Limitations and what would falsify this}

\begin{itemize}\itemsep2pt
\item Every task here is a dense statevector simulation with binomial shot sampling. The mechanisms
      are physical (control-line rectification, staggered-field mirror parity) but the data are not
      from hardware, and no device claim is made.
\item The policy menu is finite and fixed before calibration. Policies generated from calibration
      losses would require the post-selection accounting of~\citet{chen2026release}, which this paper
      does not develop; our certificate would not cover them.
\item The certification frontier depends strongly on the chosen inequality --- by a factor of four
      in $n$ between the two bounds we tried --- and the direction of the preparations-versus-shots
      tradeoff at matched budget is not robust to that choice (\S\ref{sec:withdraw}). The
      bound-independent content is the variance decomposition, not a frontier location.
\item On the independent-evaluation task the two target variants ($a=0$ and $a=1$) give identical
      risk to five decimals, i.e.\ that task's response is nearly even in $\Delta$ and does not
      exercise the odd-content axis at all. It is a transfer check for the decomposition, not for the
      sign-flip phenomenon.
\item The independent-evaluation task is small ($n=7$) and its rectification gap is an order of
      magnitude smaller than the primary task's, so it is a transfer check, not a demonstration.
\item If a hardware or publicly measured dataset showed the opposite sign for the benefit of orbit
      averaging at matched shots, the decomposition would survive but the operational reading would
      not; that is the experiment that would most directly refute the useful part of this paper.
\end{itemize}

\section*{Reproducibility}

Code: \texttt{shotcert/} (\texttt{core.py}, \texttt{qsim.py}, \texttt{tasks.py}, \texttt{physics.py},
\texttt{certify.py}, \texttt{experiments/}). Tests: \texttt{python3 -m unittest discover -s shotcert
-t . -p 'test\_*.py'} (16 tests, including the decomposition identity, the feasibility properties,
and non-degeneracy checks for every task). Results: \texttt{shotcert/results/\{control,physics,
variance\}.json}, compiled into \texttt{shotcert/RESULTS.md} by
\texttt{shotcert/compile\_results.py}; every number quoted here is emitted by that compiler.
Protocol and novelty gate: \texttt{shotcert/docs/PROTOCOL.md}, \texttt{shotcert/docs/NOVELTY\_MEMO.md}.

\bibliographystyle{plainnat}
\begin{thebibliography}{9}
\bibitem[Catak(2026)]{catak2026budget} Catak. Distribution-free generalization bounds for quantum
measurement hypotheses with separate sample and shot budgets, and a closed-form allocation rule.
arXiv:2608.24891, 2026.
\bibitem[Chen(2026a)]{chen2026release} Chen. Statistical symmetry release: certified descent by
relaxing equivariance. arXiv:2609.11470, 2026.
\bibitem[Chen(2026b)]{chen2026symmetry} Chen. Symmetry discovery: finite-measurement inference of
symmetry groups. arXiv:2610.00157, 2026.
\bibitem[Wang et~al.(2026)]{wang2026fourier} Wang, et~al. Fourier symmetrization for geometric
quantum machine learning. arXiv:2610.01874, 2026.
\bibitem[Cai et~al.(2021)]{cai2021random} Cai, et~al. Randomized symmetry in variational quantum
algorithms. arXiv:2110.13116, 2021.
\end{thebibliography}

\end{document}
